% ১৭. Practical Database Task
\section{Practical Database Task}
এই অধ্যায়ের ব্যবহারিক অংশে সাধারণত MS Access বা সমজাতীয় DBMS ব্যবহার করে ডেটাবেজ, টেবিল, কুয়েরি, ফর্ম ও রিপোর্ট তৈরি করতে বলা হয়।

\begin{practicebox}
\textbf{কাজ: StudentDB তৈরি}
\begin{enumerate}
    \item StudentDB নামে একটি ডেটাবেজ তৈরি করো।
    \item Student নামে একটি টেবিল তৈরি করো।
    \item ফিল্ড দাও: StudentID, Roll, Name, GroupName, GPA, Mobile।
    \item StudentID-কে Primary Key করো।
    \item অন্তত ৫টি রেকর্ড এন্ট্রি করো।
    \item GPA 5.00 পাওয়া ছাত্রদের দেখানোর জন্য Select query তৈরি করো।
    \item Name অনুযায়ী Ascending sort করো।
    \item Student Entry Form এবং GPA Report তৈরি করো।
\end{enumerate}
\end{practicebox}

\begin{examplebox}{নমুনা কুয়েরি}
\begin{verbatim}
SELECT Roll, Name, GPA
FROM Student
WHERE GPA = 5.00
ORDER BY Name ASC;
\end{verbatim}
\end{examplebox}

\begin{keypointbox}{ব্যবহারিক খাতার জন্য}
প্রতিটি কাজের শেষে স্ক্রিনশট/চিত্রের স্থানধারক রাখা যেতে পারে: Table Design View, Datasheet View, Query Design, Form View, Report Preview।
\end{keypointbox}
